% TÍTULO — una hoja de lectura (skill /dossier). Para quién y qué responde.
% Regenerar:  python3 graficos.py && latexmk -lualatex doc.tex
% Todo el contenido de abajo es de ejemplo: se reemplaza entero. Todo entra en una página.
\documentclass[10pt,a4paper]{article}
\usepackage[spanish,es-noshorthands]{babel}
\usepackage{dossier}
\documento{Proyecto · Título corto}{Título completo del documento}{Autor o equipo}

% Identidad del proyecto (opcional): redefinir la paleta con sus tokens, por ejemplo
% \definecolor{acento}{HTML}{22456F}

\begin{document}

\portada{Materia o proyecto · equipo · contexto}
  {Título que dice de qué trata}
  {Subtítulo: la pregunta que la hoja responde, en una línea}
  {Para quién es y fecha.}

\begin{cifras}[4]
  \cifra{Fuentes revisadas}{9}{5 controles cada una}
  \cifra{Decisiones}{12}{9 cerradas · 3 abiertas}
  \cifra{Fuera de plazo}{28,1\%}{de 3.265 tickets \fuente{R14}}
  \cifra{Filas del registro}{21.748}{31 columnas \fuente{R02}}
\end{cifras}

% Resumen: tres párrafos; con -texto, tres viñetas. El título dice el mensaje de la hoja.
\resumen[Título que dice el mensaje de la hoja]
\textbf{Qué pasó.} La respuesta primero, con sujeto y verbo, y la cifra que la sostiene
\fuente{R14}.

\textbf{Qué cambió.} Lo que el lector sabía y ya no vale, y por qué.

\textbf{Qué sigue.} La próxima decisión, con fecha.

% El pie dice primero qué se lleva el lector; después, si hace falta, cómo leer la figura.
\figura{fig/f1_ejemplo.pdf}{Quitar duplicados mueve poco la cifra; el mes de medición la mueve
mucho \fuente{R14 · R01}.}

% Dos columnas: una pieza junto al aviso, alineadas por arriba.
\noindent\begin{columna}{0.5\linewidth}
\figura{fig/f2_categorias.pdf}{2025 quedó 40 clientes por debajo de 2024 \fuente{R03}.}
\end{columna}\hfill
\begin{columna}{0.46\linewidth}
\ojo[Antes del 18/05]{Lo único que el lector no puede pasar por alto.}
\end{columna}

% Una fila por decisión: qué gana y qué cuesta, en pocas palabras.
\begin{tabularx}{\linewidth}{@{}P{3.4cm}LL@{}}
\toprule
\enc{Decisión} & \enc{Qué gana} & \enc{Qué cuesta} \\
\midrule
\textbf{Quitar duplicados} & una cifra por pedido & dos días \\
\textbf{Medir a fin de mes} & comparar meses & un corte fijo \\
\textbf{Sumar la regla nueva} & menos falsos avisos & \chip{propuesta} \\
\bottomrule
\end{tabularx}
\medskip

\begin{bloque}[Azul oscuro: hecho. Celeste: en curso. Gris: lo que falta.]
\proceso[2]{Relevamiento/A1/hecho, Datos y calidad/A2/hecho, Modelo/A3/actual, Plan y piloto/A4/futuro}
\end{bloque}

\end{document}
